% sections/introduction.tex --- Introduction
\section{Introduction}
\label{sec:introduction}

\subsection{The benchmark gap in machine-learning for audit}
\label{sec:introduction:gap}

Machine-learning approaches to accounting and audit analytics have
matured rapidly.  Deep autoencoder networks were shown to surface
journal-entry outliers at scale~\citep{schreyer2017detection};
adversarial training was extended to detect synthetic
``accounting-deepfake'' postings~\citep{schreyer2019adversarial};
gradient-boosted trees, density-based outlier detectors, and---more
recently---graph neural networks have all been reported as effective on
particular fraud and error schemes
\citep{west2016intelligent,hilal2022financial}.  In parallel, the
auditing standards themselves have moved towards explicit endorsement
of analytical procedures~\citep{vasarhelyi2004principles,isa240} and a
standardised data-exchange schema~\citep{iso21378}, both of which
presuppose that practitioners can compare detection methods on
comparable data.

Yet the empirical literature on accounting-anomaly detection has, to a
remarkable degree, advanced \emph{without a shared benchmark}.  Almost
every published evaluation operates on a private dataset---typically a
single client's general ledger, held under audit-confidentiality
constraints---and reports figures that are reproducible only by the
authors and by readers with access to the same proprietary system.
Two effects follow.  First, methodological comparison across papers is
not possible: a graph attention network reported as ``state of the
art'' on one corporate ledger may underperform a simple rule-based
detector on another, and the field has no way of detecting this from
the published numbers alone.  Second, the gap between the empirical
literature and what audit practitioners actually deploy remains
opaque: there is no widely-accepted measuring stick against which a
new technique can claim concrete improvement.

By contrast, neighbouring fields have made disciplined use of public
benchmarks for two decades.  ImageNet and CIFAR
\citep{deng2009imagenet,krizhevsky2009cifar} drove a generation of
progress in computer vision precisely because results could be
reproduced and compared at near-zero friction.  The Open Graph
Benchmark~\citep{hu2020ogb} and follow-on benchmarks for graph anomaly
detection~\citep{liu2024gadbench,tang2023bond} have begun to provide
the same service for graph machine learning.  No equivalent reference
exists in the accounting domain.

The obstacle is not principally technical: it is institutional.  Real
general ledgers cannot be released without exposing identifiable
business activity, customer and vendor relationships, and personnel
data.  Even after anonymisation, joins against public records can
re-identify entities; even after aggregation, the resulting data may
no longer support the structural analyses that audit techniques
depend on.  The community has therefore been left with a choice
between empirical realism (private, irreproducible) and
reproducibility (public, but typically toy-scale and unrealistic).

This paper takes the position that the choice is false---provided the
synthetic side is held to a high empirical and statistical standard.
We construct \bench{}, a synthetic reference benchmark, on top of the
\system{} generator~\citep{ivertowski2026datasynth} and an empirical
study of $364$ million journal-entry postings drawn from $155$ client
general ledgers~\citep{ivertowski2024apn}.  The generator is open
source~\citep{datasynth-repo}, the configuration that produces the
benchmark dataset is shipped in the same repository, and the
edge-list export, anomaly labels, and baseline implementations are
all reproducible end-to-end from a single command.

\subsection{Importance of analysing journal entries}
\label{sec:introduction:je-importance}

Journal entries are the elementary units of accounting truth.  Every
financial event in an enterprise---a sale, a goods receipt, a payroll
run, a depreciation calculation, an intercompany allocation---is
recorded as a journal entry consisting of two or more line items, the
debits of which must equal the credits.  Trial balances, subledger
reports, financial statements, and consolidated group accounts are all
derived structures: they are aggregations of, or projections from, the
underlying journal-entry table.  An auditor reasoning about an
anomaly in any of these derived structures must, at the limit, trace
the anomaly back to the journal entries that produced it.

This primacy of the journal entry is reflected in the audit standards.
ISA~240~\citep{isa240} explicitly requires the auditor to test
journal entries for indicators of management fraud, and lists
characteristics---unusual posting times, round-dollar amounts, manual
postings to typically-automated accounts, period-end clustering---that
have become standard targets for analytical procedures.  ISO~21378
\citep{iso21378} standardises the schema in which journal entries are
exported for audit data collection precisely because audit work
products converge on this layer.  Most recurring fraud typologies
documented in the practitioner literature---split transactions, ghost
employees, duplicate payments, period-cutoff manipulation, fictitious
revenue---manifest first as patterns in journal entries.

The consequence for machine learning is that the journal-entry table
is the right object of study.  Aggregated views (trial balances,
roll-up dashboards) discard the structure on which detection depends;
narrower views (single-account histories) discard the cross-account
relationships that distinguish, for example, a legitimate revenue
posting from a fictitious one.  A benchmark designed to drive
methodological progress must therefore expose journal entries in their
native granularity, with their full source-document and
inter-account context.

\subsection{Benefits of accounting-network analysis}
\label{sec:introduction:network-benefits}

Journal entries do not merely admit a network interpretation; they
\emph{are} a directed graph by construction.  This is not a new
observation: \citet{ijiri1965incidence} formalised the algebraic
relationship between double-entry bookkeeping and the
incidence matrix of a directed graph in $1965$, and the same author
later argued in his programmatic essay ``The beauty of double-entry
bookkeeping''~\citep{ijiri1993beauty} that the directed-graph
structure---accounts as nodes, postings as directed edges from
credit accounts to debit accounts---is the mathematical signature of
double-entry as an information system, and the property that
distinguishes it from single-entry record-keeping.  Subsequent work
in the audit-analytics literature
\citep{liang2021bookkeeping,guo2022picture,boersma2024audit} has
built on this foundation, treating the accounting graph as the
primary object of study rather than the journal-entry table from
which it is derived; the present author's earlier empirical
study~\citep{ivertowski2024apn} reports the structural regularities
of this graph at scale across $155$ general ledgers.

The consequence is that an accounting network is not a derived view
of the underlying postings: the postings \emph{are} edges, the
accounts \emph{are} nodes, and the journal-entry table is a
particular tabular serialisation of an object that is intrinsically
graph-shaped.  The accounting network therefore carries information
that is not present in any single row of the row-level table.

First, network structure exposes recurring \emph{motifs}.  Routine
operating cycles---procure-to-pay, order-to-cash, hire-to-retire---trace
out characteristic edge patterns whose deviations are themselves
diagnostic.  A purchase order that triggers a goods receipt that
triggers a vendor invoice that triggers a payment yields a
predictable sequence of edges across well-defined account pairs;
schemes such as duplicate payments, three-way-match overrides, and
cutoff manipulation appear as distinctive perturbations of that
sequence.

Second, network structure is the natural input for graph neural
networks (GNNs).  GNNs aggregate information across local
neighbourhoods of the graph and have been shown to outperform
non-relational baselines on a wide variety of structured-data tasks
\citep{kipf2017semi,hamilton2017inductive,velickovic2018graph,xu2019gin}.
On accounting graphs in particular, edge-level classification of
suspect postings can exploit the company a posting keeps---the other
edges incident to the same accounts, the same documents, or the same
counterparties---in ways that flat tabular models cannot.  GNN-based
graph anomaly detection has been actively benchmarked on synthetic and
public-domain graph data~\citep{liu2024gadbench,tang2023bond}, but
\emph{not} on accounting graphs of realistic size and structure.

Third, network analysis supports a class of statistical reasoning that
is otherwise hard to formalise.  Benford's law~\citep{nigrini2012benfords}
characterises distributional properties of amounts, but the same
postings additionally carry topological signatures---degree
distributions, density of strongly-connected components, betweenness
of central accounts---that fraud and error patterns alter.  A reference
benchmark that exposes both the row-level table and the induced
network-edge list is therefore strictly more useful than one that
exposes only the former.

\subsection{Contributions}
\label{sec:introduction:contributions}

This paper makes the following contributions, organised in the order in
which they appear in the remainder of the document.

\begin{enumerate}

  \item \textbf{Empirical foundation.}  We summarise---and revisit
    where appropriate---the empirical study of journal-entry
    line-item distributions on a corpus of $364$ million postings
    drawn from $155$ client general ledgers \citep{ivertowski2024apn}.
    The recovered distribution informs the parameters that
    \system{}~\citep{ivertowski2026datasynth} uses to generate
    accounting networks of realistic structure
    (Section~\ref{sec:empirical}).

  \item \textbf{Network construction methods.}  We present, with full
    mathematical statement, the five Cartesian-product methods
    (A through E) by which a balanced journal entry can be
    materialised as a set of debit\,$\to$\,credit edges, and show
    their correspondence to the standard
    financial-statement-consolidation operation
    (Section~\ref{sec:network}).

  \item \textbf{Reference graph schema.}  We define a $13$-column flat
    edge-list export (\texttt{graphs/je\_network.csv}) with
    deterministic UUID v5 edge identifiers, joinable column-for-column
    against the journal-entry table and against the document-flow
    chain via a \texttt{predecessor\_line\_id} field
    (Section~\ref{sec:graph}).  The schema is sufficient to reproduce
    all baselines reported here without additional pre-processing.

  \item \textbf{Benchmark design.}  We open-source \bench{}: a
    fully-specified benchmark with three reproducible splits
    (temporal, entity, fraud-typology), a labelling discipline that
    surfaces over $130$ anomaly subtypes at the edge level, and a
    leaderboard protocol that includes a leakage audit
    (Section~\ref{sec:benchmark}).

  \item \textbf{Strong baselines and initial evaluation.}  We provide
    a baseline suite covering five GNN architectures, gradient-boosted
    trees, two classical anomaly detectors, and a rule-based
    reference, all trained from a single declarative configuration,
    and report initial numbers on the released splits
    (Section~\ref{sec:evaluation}).

\end{enumerate}

\subsection{Paper organisation}
\label{sec:introduction:organisation}

Section~\ref{sec:related-work} surveys the prior literature on
accounting-anomaly detection, graph anomaly benchmarks, and
synthetic-data generation for finance.  Section~\ref{sec:empirical}
revisits the empirical foundations.  Section~\ref{sec:datasynth}
summarises the \system{} generator that produces the
benchmark dataset.  Section~\ref{sec:network} states the network
construction methods.  Section~\ref{sec:graph} describes the
reference graph schema.  Section~\ref{sec:benchmark} fixes the
benchmark protocol.  Section~\ref{sec:evaluation} reports baseline
results.  Section~\ref{sec:discussion} interprets them and discusses
threats to validity, and Section~\ref{sec:future-work} lists open
problems.  Section~\ref{sec:conclusion} concludes.
